% GENERATED FILE. Edit canonical lessons, then run npm run generate:books.
% canonical-source-hash: fnv1a64:a7f361c41b86c201
% canonical-lessons: RU-C129-soft, RU-C129-hard, RU-C129-round, RU-C129-wide, RU-C129-narrow

\chapter{Soft, Hard, Round, Wide, Narrow}
\label{ch:ru-soft-hard-round-wide-narrow}
\hlchaptermodality{129}

\begin{chapteropening}
\textbf{By the end of this chapter:} \emph{I can say soft, hard, round, wide and narrow.}

Complete the last lesson of chapter 129: I can say soft, hard, round, wide and narrow.
\end{chapteropening}

\section[\texorpdfstring{\ru{мягкий} (myákhkiy)}{мягкий (myákhkiy)}]{\ru{мягкий} (myákhkiy) — soft}

\label{lesson:RU-C129-soft}



\begin{quote}
Before the new one: say the Russian for wet, then the Russian for dry.
\end{quote}

\subsection*{You'll want to know: \ru{мягкий}}
\textbf{\ru{мягкий}} (\emph{myákhkiy}) — \textquotedblleft{}soft\textquotedblright{}.

\begin{cousinweb}
\textbf{\ru{мягкий}} $\leftarrow$ Proto-Slavic \emph{*mękŭkŭ} — inherited. Vasmer relates it to Lithuanian \emph{mìnkštas}, \textquotedblleft{}soft\textquotedblright{}, and \emph{mìnkyti}, \textquotedblleft{}to knead\textquotedblright{}, and Greek \emph{mássō}, \textquotedblleft{}I knead\textquotedblright{}; \emph{\ru{мука}}, \textquotedblleft{}flour\textquotedblright{}, shares its root.
\end{cousinweb}

\begin{grammarlens}[title={What you've built}]
One more word for this chapter.
\end{grammarlens}

\subsection*{Your turn}
\begin{itemize}
\raggedright
  \item \emph{Say it:} \emph{myákhkiy}
  \item \emph{Say it:} \emph{myákhkiy}, once more
  \item \emph{Say it:} say \emph{myákhkiy}
  \item \emph{Recall:} say the Russian for wet, then the Russian for dry, then say \emph{myákhkiy} again
\end{itemize}

\begin{tcolorbox}[breakable,colback=teal!4,colframe=teal!35!black,title={Before you move on}]
What does \ru{мягкий} mean? (\textquotedblleft{}Soft\textquotedblright{}.) Say it once more.
\end{tcolorbox}
\section[\texorpdfstring{\ru{твёрдый} (tvyórdyy)}{твёрдый (tvyórdyy)}]{\ru{твёрдый} (tvyórdyy) — hard}

\label{lesson:RU-C129-hard}



\begin{quote}
Before the new one: say the Russian for dry, then the Russian for soft.
\end{quote}

\subsection*{You'll want to know: \ru{твёрдый}}
\textbf{\ru{твёрдый}} (\emph{tvyórdyy}) — \textquotedblleft{}hard\textquotedblright{}.

\begin{cousinweb}
\textbf{\ru{твёрдый}} $\leftarrow$ Proto-Slavic \emph{*tvĭrdŭ} — inherited, a cousin of Lithuanian \emph{tvìrtas}, \textquotedblleft{}firm\textquotedblright{}, and Latvian \emph{tvirts}, \textquotedblleft{}strong\textquotedblright{} (Vasmer); \emph{\ru{творить}}, \textquotedblleft{}to create\textquotedblright{}, shares its root, Vasmer says.
\end{cousinweb}

\begin{grammarlens}[title={What you've built}]
One more word for this chapter.
\end{grammarlens}

\subsection*{Your turn}
\begin{itemize}
\raggedright
  \item \emph{Say it:} \emph{tvyórdyy}
  \item \emph{Say it:} \emph{tvyórdyy}, once more
  \item \emph{Say it:} say \emph{tvyórdyy}
  \item \emph{Recall:} say the Russian for dry, then the Russian for soft, then say \emph{tvyórdyy} again
\end{itemize}

\begin{tcolorbox}[breakable,colback=teal!4,colframe=teal!35!black,title={Before you move on}]
What does \ru{твёрдый} mean? (\textquotedblleft{}Hard\textquotedblright{}.) Say it once more.
\end{tcolorbox}
\section[\texorpdfstring{\ru{круглый} (krúglyy)}{круглый (krúglyy)}]{\ru{круглый} (krúglyy) — round}

\label{lesson:RU-C129-round}



\begin{quote}
Before the new one: say the Russian for soft, then the Russian for hard.
\end{quote}

\subsection*{You'll want to know: \ru{круглый}}
\textbf{\ru{круглый}} (\emph{krúglyy}) — \textquotedblleft{}round\textquotedblright{}.

\begin{cousinweb}
\textbf{\ru{круглый}} is built on \emph{\ru{круг}}, \textquotedblleft{}a circle, a ring\textquotedblright{} (Vasmer), with the adjective ending \emph{-\ru{лый}}.
\end{cousinweb}

\begin{grammarlens}[title={What you've built}]
One more word for this chapter.
\end{grammarlens}

\subsection*{Your turn}
\begin{itemize}
\raggedright
  \item \emph{Say it:} \emph{krúglyy}
  \item \emph{Say it:} \emph{krúglyy}, once more
  \item \emph{Say it:} say \emph{krúglyy}
  \item \emph{Recall:} say the Russian for soft, then the Russian for hard, then say \emph{krúglyy} again
\end{itemize}

\begin{tcolorbox}[breakable,colback=teal!4,colframe=teal!35!black,title={Before you move on}]
What does \ru{круглый} mean? (\textquotedblleft{}Round\textquotedblright{}.) Say it once more.
\end{tcolorbox}
\section[\texorpdfstring{\ru{широкий} (shirókiy)}{широкий (shirókiy)}]{\ru{широкий} (shirókiy) — wide}

\label{lesson:RU-C129-wide}



\begin{quote}
Before the new one: say the Russian for hard, then the Russian for round.
\end{quote}

\subsection*{You'll want to know: \ru{широкий}}
\textbf{\ru{широкий}} (\emph{shirókiy}) — \textquotedblleft{}wide\textquotedblright{}.

\begin{cousinweb}
\textbf{\ru{широкий}} $\leftarrow$ Proto-Slavic \emph{*širokŭ} — inherited, from \emph{*širŭ}, \textquotedblleft{}wide\textquotedblright{}. Vasmer calls its further links unclear; a comparison with Gothic \emph{skeirs}, \textquotedblleft{}clear\textquotedblright{}, is hard to square with the sounds.
\end{cousinweb}

\begin{grammarlens}[title={What you've built}]
One more word for this chapter.
\end{grammarlens}

\subsection*{Your turn}
\begin{itemize}
\raggedright
  \item \emph{Say it:} \emph{shirókiy}
  \item \emph{Say it:} \emph{shirókiy}, once more
  \item \emph{Say it:} say \emph{shirókiy}
  \item \emph{Recall:} say the Russian for hard, then the Russian for round, then say \emph{shirókiy} again
\end{itemize}

\begin{tcolorbox}[breakable,colback=teal!4,colframe=teal!35!black,title={Before you move on}]
What does \ru{широкий} mean? (\textquotedblleft{}Wide\textquotedblright{}.) Say it once more.
\end{tcolorbox}
\section[\texorpdfstring{\ru{узкий} (úzkiy)}{узкий (úzkiy)}]{\ru{узкий} (úzkiy) — narrow}

\label{lesson:RU-C129-narrow}



\begin{quote}
Before the new one: say the Russian for round, then the Russian for wide.
\end{quote}

\subsection*{You'll want to know: \ru{узкий}}
\textbf{\ru{узкий}} (\emph{úzkiy}) — \textquotedblleft{}narrow\textquotedblright{}.

\begin{cousinweb}
\textbf{\ru{узкий}} $\leftarrow$ Proto-Slavic \emph{*ǫzŭkŭ} — inherited. Vasmer lists Sanskrit \emph{aṃhú-}, \textquotedblleft{}narrow\textquotedblright{}, Gothic \emph{aggwus}, Lithuanian \emph{añkštas} and Latin \emph{angustus}, all \textquotedblleft{}narrow\textquotedblright{}, and Greek \emph{ánchō}, \textquotedblleft{}I squeeze\textquotedblright{}.
\end{cousinweb}

\begin{grammarlens}[title={What you've built}]
One more word for this chapter.
\end{grammarlens}

\subsection*{Your turn}
\begin{itemize}
\raggedright
  \item \emph{Say it:} \emph{úzkiy}
  \item \emph{Say it:} \emph{úzkiy}, once more
  \item \emph{Say it:} say \emph{úzkiy}
  \item \emph{Recall:} say the Russian for round, then the Russian for wide, then say \emph{úzkiy} again
\end{itemize}

\begin{tcolorbox}[breakable,colback=teal!4,colframe=teal!35!black,title={Before you move on}]
What does \ru{узкий} mean? (\textquotedblleft{}Narrow\textquotedblright{}.) Say it once more.
\end{tcolorbox}
